\section{A learning roadmap: what to study next}
\label{sec:roadmap}

A suggested order to go from ``I read this guide'' to ``I can change the code.''

\subsection*{Day 1 --- orient and run}
\begin{enumerate}[leftmargin=1.6em,itemsep=3pt]
  \item Re-read \autoref{sec:bigpicture} and
    \autoref{sec:physics}.1--\ref{sec:physics}.3 until ``heat $\ne$ temperature''
    and the heat equation feel natural.
  \item Do the setup (\autoref{sec:running}.1), run the tests
    (\autoref{sec:running}.2), verify the headline numbers
    (\autoref{sec:running}.3).
  \item Run \file{scripts/analysis/phase2\_depth\_convergence.py}
    (\autoref{sec:running}.5) and stare at the output figure while re-reading
    \autoref{sec:skin}. This makes skin depth and the geothermal gradient
    \emph{concrete}.
\end{enumerate}

\subsection*{Day 2 --- read the engine}
\begin{enumerate}[leftmargin=1.6em,itemsep=3pt,start=4]
  \item Read \file{lunar/constants.py} then \file{lunar/config.py} end to end ---
    they are short and define ``the experiment.''
  \item Read \file{lunar/properties.py} (\code{conductivity\_hayne},
    \code{density\_hayne}, \code{specific\_heat}) alongside
    \autoref{sec:properties}--\ref{sec:conductivity}.
  \item Read \file{lunar/grid.py} alongside \autoref{sec:grid}.
  \item Skim \file{lunar/solver.py}: focus on the \emph{docstring},
    \code{surface\_energy\_balance\_residual}, and \code{\_thomas}. Don't try to
    absorb every index in \code{\_step} yet.
\end{enumerate}

\subsection*{Day 3 --- the centerpiece}
\begin{enumerate}[leftmargin=1.6em,itemsep=3pt,start=8]
  \item Read the full docstring of \file{lunar/equilibrium.py} (it is a teaching
    text in itself), then \code{solve\_periodic\_equilibrium}, with
    \autoref{sec:anchor} open beside it.
  \item Read \file{tests/test\_equilibrium.py} --- the tests \emph{are} the
    precise statement of what ``correct'' means here.
  \item Optional: run \code{equilibrium\_demo.ipynb} to watch brute-force
    converge onto the fast method.
\end{enumerate}

\subsection*{Day 4 --- the science}
\begin{enumerate}[leftmargin=1.6em,itemsep=3pt,start=11]
  \item Read \file{scripts/pipeline/retrieve\_kd.py} top to bottom with
    \autoref{sec:retrieval} open. Follow one \Kd{} value through \code{run\_with}
    $\to$ \code{run\_kd\_sweep\_extended} $\to$ \code{kd\_star\_from\_residuals}.
  \item Read \file{lunar/apollo\_helpers.py} to see how raw data becomes
    $T_{\mathrm{eq}}$.
  \item Skim the paper abstract, plain-language summary, and \S2 (methods) in
    \file{paper/letter/letter.tex}.
\end{enumerate}

\subsection*{Then --- make a change (the real test of understanding)}
\begin{enumerate}[leftmargin=1.6em,itemsep=3pt,start=14]
  \item Try a tiny experiment: in a scratch script, call
    \code{run\_with(SITES['A15'], kd=...)} at a couple of \Kd{} values and print
    the temperature at $1.4$~m. Watch it move toward the data as \Kd{} approaches
    \num{4.58}.
  \item Read \file{docs/FLAG\_REPORT.md} to see how the project audits
    \emph{itself} --- this is how careful computational science is actually done.
\end{enumerate}

\subsection*{Deeper study, when ready}
\begin{itemize}[leftmargin=1.4em,itemsep=3pt]
  \item \textbf{Numerical methods:} the finite-volume method, implicit vs
    explicit schemes, stability (any intro numerical-PDE text).
  \item \textbf{Heat transfer:} Fourier's law and the 1-D heat equation (any
    intro heat-transfer text).
  \item \textbf{Python scientific stack:} NumPy basics (arrays, broadcasting,
    \code{np.interp}, \code{np.gradient}) --- used everywhere here.
  \item \textbf{Statistics:} the bootstrap and confidence intervals.
\end{itemize}

\begin{keyidea}
The single most effective habit: keep this guide open on one side and the actual
code file on the other. Every concept here points at a real function you can
open, run, and modify. Understanding grows fastest when the words and the code
are side by side.
\end{keyidea}
